\section{Results}
\label{sec:results}

\subsection{Data harmonisation and SPARK-only target}

The two sumstats files are on different genome builds (iPSYCH--PGC on
GRCh37, the Stein-Lab meta-analysis on GRCh38). A na\"{i}ve
\texttt{chr:pos} inner-join therefore matched only $\sim 73$\,k of the
$\sim 9$\,M SNPs in either file; the rsID-based join described in
Section~\ref{sec:method} recovered $7\,413\,400$ shared SNPs, of which
$6\,309\,737$ survived QC (37\,101 in MHC). The harmonised file has median absolute discovery Z about 0.72 and
median absolute meta Z about 0.72 as well, consistent with mild
polygenic inflation.

We checked the SPARK-only back-calculation by plugging the recovered
values back into the same blending rule as the meta-analysis: the
reconstructed meta Z matches the published file to machine precision
(maximum error about $10^{-15}$, i.e.\ rounding only). The derived
SPARK-only Z-scores have median absolute value about 0.72 and barely
correlate with the discovery Z (correlation 0.035), as we expect when
the independent SPARK signal is weak and noisy at each SNP.

The Manhattan and QQ plots of the discovery sumstats
(Fig.~\ref{fig:manhattan},~\ref{fig:qq}) reproduce the genome-wide
significant peaks reported by Grove et al.~\cite{grove2019} and exhibit
the characteristic upward-bending QQ tail consistent with polygenic
inflation rather than confounding.

\begin{figure}[h]
  \centering
  \includegraphics[width=\linewidth]{fig_manhattan.pdf}
  \caption{Manhattan plot of the iPSYCH--PGC discovery sumstats
    (6\,309\,737 SNPs after harmonisation and QC).
    The horizontal red line marks $p = 5 \times 10^{-8}$.}
  \label{fig:manhattan}
\end{figure}

\begin{figure}[h]
  \centering
  \includegraphics[width=0.65\linewidth]{fig_qq.pdf}
  \caption{QQ plot of significance ($-\log_{10}$ of the $p$-value) for
    the iPSYCH--PGC discovery sumstats. The diagonal line is the null
    expectation.}
  \label{fig:qq}
\end{figure}

\subsection{Held-out predictive accuracy}

We evaluate every method on chr~21--22 (170\,039 test SNPs after
removing MHC). Table~\ref{tab:eval} reports how well each method's
predicted Z-score tracks the SPARK-only Z (Pearson correlation $r$ with
bootstrapped 95\,\% intervals from 1000 resamples of SNPs), the AUC for
flagging large absolute SPARK Z, and the calibration slope.
Fig.~\ref{fig:bars} renders the same numbers visually, and
Fig.~\ref{fig:scatter} shows the per-method predicted-vs-observed
scatter on the same hold-out.

\begin{table}[h]
  \centering
  \small
  \setlength{\tabcolsep}{4pt}
  \begin{tabular}{lrrr}
    \toprule
    Method & $n$ & $r$ [95\,\% CI] & AUC \\
    \midrule
    \textbf{ml\_xgb}     & 170\,039 & \textbf{0.058} [0.053, 0.063] & 0.488 \\
    ml\_mlp              & 170\,039 & 0.048 [0.043, 0.053] & 0.499 \\
    PRS-CS-EB            & 170\,039 & 0.041 [0.036, 0.046] & 0.499 \\
    GNN (no graph)       & 170\,039 & 0.037 [0.033, 0.042] & 0.499 \\
    \midrule
    P+T $p \!\leq\! 1.0$     &      213 & $-0.046$ [$-0.18$, $0.09$] & 0.587 \\
    P+T $p \!\leq\! 5e\!-\!2$ & 191 & $-0.045$ [$-0.17$, $0.09$] & 0.537 \\
    P+T $p \!\leq\! 10^{-3}$ &       37 & $-0.097$ [$-0.36$, $0.20$] & 0.578 \\
    P+T $p \!\leq\! 10^{-5}$ &        2 & ---                          & --- \\
    P+T $p \!\leq\! 5e\!-\!8$ &  0 & ---                          & --- \\
    \bottomrule
  \end{tabular}
  \caption{Held-out test-fold metrics (chr~21--22, MHC excluded). Rows
    above the inner rule are evaluated on every test SNP; P+T rows
    include only the SNPs that survived clumping at the listed
    threshold (\texttt{---} where $n \!<\! 3$ makes the metric
    meaningless).}
  \label{tab:eval}
\end{table}

\begin{figure}[h]
  \centering
  \includegraphics[width=\linewidth]{fig_method_bars.pdf}
  \caption{Held-out predictive accuracy. Left: correlation between
    predicted and SPARK-only Z on chr~21--22, with 95\,\% CIs. Right:
    AUC for detecting large absolute SPARK-only Z.}
  \label{fig:bars}
\end{figure}

\begin{figure}[h]
  \centering
  \includegraphics[width=\linewidth]{fig_method_scatter.pdf}
  \caption{Per-method predicted (x-axis) versus observed SPARK-only
    Z-score (y-axis) on chr~21--22, hex-binned. Each in-text $r$ value
    is annotated.}
  \label{fig:scatter}
\end{figure}

The four whole-genome methods produce a clear ordering. The
\textbf{XGBoost re-weighter} attains the highest Pearson $r = 0.058$
(95\,\% CI $[0.053, 0.063]$), beating the empirical-Bayes baseline by
0.017 and the discovery-only \textit{GNN-no-graph} fallback by 0.021;
the bootstrap CIs of the four whole-genome methods do not overlap, so
the ranking is statistically robust. The \textbf{MLP re-weighter} comes
second ($r = 0.048$). Both ML methods improve over the empirical-Bayes
baseline despite using only the discovery Z, $|Z|$, INFO score, and
LD-window summaries -- the gene-level functional annotations sketched
in Section~\ref{sec:method} were not available in this run because the
required reference files (Gencode, GTEx, ENCODE, SFARI Gene, STRING)
are gated by separate download steps; we expect the gain to widen
materially once those features are added.

The \textbf{P+T baselines} are dominated on this evaluation. Even the
most permissive threshold ($p \leq 1$) retains only 213 chr~21--22
SNPs after clumping, and at $p \leq 5\!\times\!10^{-8}$ no chr~21--22
SNP survives -- a direct illustration of P+T's known limitation for a
highly polygenic trait. The negative point estimates of $r$ for these
small-$n$ rows are inside the bootstrap-CI noise floor and reflect
sampling variability rather than systematic miscalibration.

The AUC for the ``large absolute SPARK Z'' task hovers near 0.5 for
the whole-genome methods, which is expected: on chromosomes 21--22,
very few SNPs show a strong independent SPARK signal at the effective
sample size of the SPARK-only slice (about 12\,k people), so this
metric has little power here and should not be over-interpreted.

\subsection{Calibration}

Calibration (plotting observed SPARK-only Z against predicted Z on the
test chromosomes) tells a complementary story. The MLP is closest to
a one-to-one line (slope about 0.9). XGBoost predicts somewhat too
strongly (slope about 1.8). The empirical-Bayes baseline is far too
timid (slope about 0.06 -- its predicted Z values are much smaller than
the observed ones). In short, the ML models that are trained directly
on the held-out target behave more sensibly than the fixed shrinkage
rule.

\subsection{Stratified analyses}

The supplementary file \texttt{results/tables/eval\_strata.tsv} splits
results by imputation quality (INFO). The whole-genome methods all do
best when INFO is above 0.95 (correlations roughly 0.06--0.08), and
the relative ordering of methods is preserved. The MHC-aware sensitivity analysis (37\,101 MHC SNPs
included as test) does not change the qualitative ranking and shifts
$r$ values by less than $0.005$.

\subsection{Biological enrichment}

The enrichment analysis described in Section~\ref{sec:method} is gated
by the optional reference annotation files. Because the Gencode~v44
gene BED, SFARI Gene catalogue, STRING~v12 edges, and GTEx eQTL tracks
were not available locally for this run, the SNP-to-gene mapping was
empty and the per-method gene-set Fisher tests in
\texttt{enrichment.tsv} therefore record a \texttt{missing-ref} status
rather than $p$-values. The pipeline as published (\texttt{make
reference} in \texttt{code/Makefile}) lists the exact downloads
required to populate this table.
